\chapter{Conclusions}

The project was overall a success. The analysis tool built offers a direct performance improvement over the classical 
techniques for conducting ILP limit studies. All the minimum success criteria were met, and nearly all the extensions have 
been implemented with varying levels of success. The associativity of the matrix approach has been shown to offer 
clear pathways for both data-driven and parallel optimisations.

\section{Lessons learned}



\section{Future work}

While the core of the project is complete, there are further improvements that could be made. Listed are some shortcomings 
that could be addressed in the future.

\begin{itemize}
  \item The GPU pipeline did not make full use of the compute available. More work could be done to optimise the memory 
    layout for contiguous storage, implement a double-buffered dispatch architecture to hide overhead, and make use of
    a hybrid system that simultaneously multiplies matrices on CPU threads to maximise throughput.
  \item The memory usage of the tool is technically unbounded, and it would often crash during benchmarks, or even stall 
    due to spillover into swap space. More work could be done to ensure that the maximum memory usage is safely limited 
    during analysis of realistic traces of trillions of dynamic instructions.
  \item I/O and parsing forms a sequential bottleneck for the analysis tool when loading the initial expression. Work 
    could be done to implement an alternative binary-based format to speed up loading and reduce the bottleneck.
  \item The analysis tool only tracked register dependencies. While this is still useful for smaller subtraces that keep their 
    working set in registers, programs will often spill registers to memory. This happened with all the sample traces,
    and so reported ILP limits were higher than actual values. More work could be done to assess the impact of the data-driven 
    optimisations on the more complex structure of a matrix expression that tracks memory dependencies.
\end{itemize}

Overall, the project achieved its goals, and I think it was an interesting project, bridging gaps between concepts in 
processor design, graph theory, pattern exploitation and parallel computing.
